\documentclass[12pt]{article}
\usepackage{amsmath,amssymb,amsthm}
\usepackage{mathtools}
\usepackage{hyperref}
\usepackage{cleveref}
\usepackage{graphicx}
\usepackage{natbib}
\usepackage[margin=1in]{geometry}

\newtheorem{theorem}{Theorem}[section]
\newtheorem{proposition}[theorem]{Proposition}
\newtheorem{corollary}[theorem]{Corollary}
\newtheorem{lemma}[theorem]{Lemma}
\newtheorem{remark}[theorem]{Remark}
\newtheorem{definition}[theorem]{Definition}
\newtheorem{assumption}{Assumption}

\newcommand{\R}{\mathbb{R}}
\newcommand{\E}{\mathbb{E}}
\newcommand{\Q}{\mathbb{Q}}
\newcommand{\dd}{\mathrm{d}}
\newcommand{\vep}{\varepsilon}

\title{Solving the Correlated Variational Problem for Short-Maturity Variance Options in Local-Stochastic Volatility Models}
\author{}
\date{\today}

\begin{document}
\maketitle

\begin{abstract}
We study the short-maturity asymptotics for options on realized variance in local-stochastic
volatility (LSV) models with non-zero correlation $\rho$ between asset and volatility noises.
Pirjol, Wang, and Zhu~(2025) showed that the leading-order asymptotics are governed by
a two-dimensional variational problem, which they solved explicitly for $\rho=0$ and
$\rho=\pm 1$, but left open for general $\rho\in(-1,1)\setminus\{0\}$. We develop three
complementary approaches to this open problem under constant volatility-of-volatility
$\sigma(v)\equiv\sigma_0$. First, we establish that the Euler--Lagrange system admits a
Hamiltonian first integral, providing an energy conservation law that constrains the
optimal trajectories. Second, we introduce a one-parameter family of upper bounds,
indexed by a parameter $\alpha\in\R$, that reduce the two-dimensional problem to a
one-dimensional optimization involving the Asian option rate function with modified local
volatility. The optimal bound within this family provides up to 75\% improvement over
the best previously known bound in the leverage regime ($\rho<0$) and is
exact when the local volatility function is constant. Third, we derive
a perturbation expansion of the rate function in powers of $\rho$ around $\rho=0$, with
explicit computable coefficients determined by the known uncorrelated optimizers. We
analyze the boundary behavior as $\rho\to\pm 1$ and as the strike approaches extreme values.
Numerical experiments in the Tanh local volatility model confirm the accuracy of our
bounds and the perturbation expansion. Our results provide usable pricing approximations
for options on the recently launched CBOE S\&P~500 Variance Futures.
\end{abstract}

\noindent\textbf{Keywords:} variance options, short-maturity asymptotics, local-stochastic
volatility, large deviations, variational problem, Hamiltonian mechanics

\noindent\textbf{MSC2020:} 91G20, 60F10, 49J15

\section{Introduction}

Options on realized variance are derivative contracts whose payoffs are linked to the
annualized realized variance of an underlying asset's return. Following the launch of the
CBOE S\&P~500 Variance Futures (Ticker:~VA) in September~2024, there is renewed
interest in analytical pricing methods for these instruments.

Under a general local-stochastic volatility (LSV) model, the asset price~$S_t$ and
variance process~$V_t$ satisfy
\begin{align}
\frac{\dd S_t}{S_t} &= \eta(S_t)\sqrt{V_t}\bigl(\rho\,\dd Z_t
  + \sqrt{1-\rho^2}\,\dd W_t\bigr), \label{eq:S} \\
\frac{\dd V_t}{V_t} &= \mu(V_t)\,\dd t + \sigma(V_t)\,\dd Z_t, \label{eq:V}
\end{align}
where $W_t,Z_t$ are independent standard Brownian motions, $\rho\in(-1,1)$ is the
correlation parameter, and $\eta(\cdot),\sigma(\cdot),\mu(\cdot)$ are the local volatility,
volatility-of-volatility, and drift functions respectively. The realized variance over
$[0,T]$ is the quadratic variation
\begin{equation}\label{eq:RV}
[{\log S}]_T = \int_0^T V_s\,\eta^2(S_s)\,\dd s.
\end{equation}
A variance call option with strike~$K$ and maturity~$T$ has price
\begin{equation}\label{eq:Cprice}
C(T) = e^{-rT}\,\E\biggl[\biggl(\frac{1}{T}\int_0^T V_s\,\eta^2(S_s)\,\dd s
  - K\biggr)^{\!+}\,\biggr].
\end{equation}

Pirjol, Wang, and Zhu~\cite{PWZ2025} established that the short-maturity asymptotics
for out-of-the-money (OTM) variance options are governed by the rate function
\begin{equation}\label{eq:Irho}
\lim_{T\to 0} T\log C(T) = -I_\rho(S_0,V_0,K),
\end{equation}
where $I_\rho$ is the solution to the two-dimensional variational problem
\begin{equation}\label{eq:varprob}
I_\rho(S_0,V_0,K) = \inf_{\substack{g(0)=\log S_0,\;h(0)=\log V_0 \\
  \int_0^1 e^{h(t)}\eta^2(e^{g(t)})\,\dd t = K}} \Lambda_\rho[g,h],
\end{equation}
with the action functional
\begin{equation}\label{eq:action}
\Lambda_\rho[g,h] = \frac{1}{2(1-\rho^2)}\int_0^1
  \biggl(\frac{g'(t)}{\eta(e^{g(t)})\sqrt{e^{h(t)}}}
  - \frac{\rho\,h'(t)}{\sigma(e^{h(t)})}\biggr)^{\!2}\dd t
  + \frac{1}{2}\int_0^1\biggl(\frac{h'(t)}{\sigma(e^{h(t)})}\biggr)^{\!2}\dd t.
\end{equation}

For $\rho=0$, the variational problem decouples and was solved explicitly by
\cite{PWZ2025}, reducing to a one-dimensional optimization involving the Asian option
rate function~$J(V_0,z)$ from~\cite{PZ2023}. For $\rho=\pm 1$, the singular coefficient
$1/(1-\rho^2)$ forces the integrand of the first term to vanish, again reducing to a
one-dimensional problem. For general $\rho\in(-1,1)\setminus\{0\}$, only upper and lower
bounds were obtained.

\paragraph{Related work.}
The connection between variance option pricing and Asian option asymptotics was
exploited in~\cite{PZ2024} for the LSV model, where a similar two-dimensional
variational problem was solved via an expansion in log-moneyness. The proportionality
ansatz technique used here is methodologically related to the coupling approach
in~\cite{PZ2024}, though our square-root coupling~\eqref{eq:Falpha} differs from the
logarithmic coupling used there. In the rough volatility setting, Lacombe, Muguruza,
and Stone~\cite{LMS2021} studied short-maturity asymptotics for variance derivatives
with non-zero correlation, deriving the implied volatility smile slope at leading order.
Although their framework (multi-factor rough volatility) differs from the Markovian LSV
model considered here, both works address the fundamental challenge of incorporating
correlation into path-dependent option asymptotics. The geometric approach to LSV
asymptotics, treating the rate function as a geodesic problem on a Riemannian manifold,
was developed by Henry-Labord\`ere~\cite{HenryLabordere} and rigorously justified by
Forde and Jacquier~\cite{FordeJacquier}.

\subsection{Our contributions}

We develop three complementary approaches to the open problem of solving~\eqref{eq:varprob}
for general~$\rho$, under the assumption of constant volatility-of-volatility
$\sigma(v)\equiv\sigma_0$.

\begin{enumerate}
\item \textbf{Hamiltonian conservation} (\Cref{thm:hamiltonian}):
  We observe that the Euler--Lagrange system for~\eqref{eq:varprob} admits a first integral
  (conserved energy), reducing the effective dimension of the phase space. This provides
  a structural constraint on all optimal trajectories.

\item \textbf{Alpha-parametric bounds} (\Cref{thm:alpha}):
  We introduce a one-parameter family of upper bounds
  \[
  I_\rho \leq \min_{\alpha\in\R}\, C(\alpha,\rho)\cdot J_\alpha(V_0,K),
  \]
  where $C(\alpha,\rho) = (\alpha^2-2\alpha\rho+1)/(1-\rho^2)$ and $J_\alpha$ is a
  one-dimensional rate function computable via the methods of~\cite{PZ2023}.
  Optimizing over~$\alpha$ yields bounds that can significantly improve upon the existing
  bounds of~\cite{PWZ2025}, with up to 75\% improvement in the leverage regime
  ($\rho<0$). The bound is exact when $\eta$ is constant (\Cref{cor:exactconst}).
  The square-root coupling~\eqref{eq:Falpha} in our ansatz is new and complements the
  logarithmic coupling of Proposition~17 in~\cite{PWZ2025}; the two approaches are
  methodologically related to the coupling technique of~\cite{PZ2024} for Asian options
  in LSV models.

\item \textbf{Perturbation expansion} (\Cref{thm:perturb}):
  We derive the Taylor expansion $I_\rho = I_0 + \rho\,I_1 + O(\rho^2)$, where the
  first-order coefficient~$I_1$ is an explicit integral involving the known $\rho=0$
  optimizers.
\end{enumerate}

The paper is organized as follows. \Cref{sec:hamiltonian} establishes the Hamiltonian
structure. \Cref{sec:alpha} develops the alpha-parametric bounds.
\Cref{sec:perturbation} derives the perturbation expansion. \Cref{sec:boundary}
analyzes boundary behavior. \Cref{sec:numerics} presents numerical experiments.


\section{Hamiltonian Structure}\label{sec:hamiltonian}

Throughout this paper, we assume $\sigma(v)\equiv\sigma_0>0$ (constant
volatility-of-volatility). We introduce the velocity variables
\begin{equation}\label{eq:ab}
a(t) = \frac{g'(t)}{\eta(e^{g(t)})\sqrt{e^{h(t)}}}, \qquad
b(t) = \frac{h'(t)}{\sigma_0}.
\end{equation}
In these variables, the action functional~\eqref{eq:action} takes the form
\begin{equation}\label{eq:action_ab}
\Lambda_\rho[g,h] = \frac{1}{2(1-\rho^2)}\int_0^1\bigl(a(t)-\rho\,b(t)\bigr)^2\,\dd t
  + \frac{1}{2}\int_0^1 b(t)^2\,\dd t.
\end{equation}
Expanding: since $(a-\rho b)^2/(1-\rho^2) + b^2 = [(a-\rho b)^2 + (1-\rho^2)b^2]/(1-\rho^2) = (a^2-2\rho ab+b^2)/(1-\rho^2)$, we obtain
\begin{equation}\label{eq:action_expanded}
\Lambda_\rho[g,h] = \frac{1}{2(1-\rho^2)}\int_0^1
  \bigl(a^2 - 2\rho\,a\,b + b^2\bigr)\,\dd t.
\end{equation}

The constrained variational problem~\eqref{eq:varprob} with Lagrange multiplier
$\lambda$ gives the augmented Lagrangian density
\begin{equation}\label{eq:Lagrangian}
\mathcal{L}(g,g',h,h') = \frac{1}{2(1-\rho^2)}\biggl(\frac{g'}{\eta(e^g)\sqrt{e^h}}
  - \frac{\rho\,h'}{\sigma_0}\biggr)^{\!2}
  + \frac{1}{2}\biggl(\frac{h'}{\sigma_0}\biggr)^{\!2}
  + \lambda\,e^h\eta^2(e^g).
\end{equation}

\begin{proposition}[Hamiltonian conservation]\label{thm:hamiltonian}
Assume $\sigma(v)\equiv\sigma_0$. Let $(g,h)$ be a smooth critical point of the
constrained variational problem~\eqref{eq:varprob}. Define the Hamiltonian
\begin{equation}\label{eq:H}
\mathcal{H}(t) = \frac{1}{2(1-\rho^2)}\bigl(a(t)^2 - 2\rho\,a(t)\,b(t) + b(t)^2\bigr)
  - \lambda\,e^{h(t)}\eta^2(e^{g(t)}).
\end{equation}
Then $\mathcal{H}(t) = E$ for all $t\in[0,1]$, where $E$ is a constant determined by
the boundary conditions.
\end{proposition}

\begin{proof}
Since $\sigma(v)\equiv\sigma_0$, the Lagrangian density~\eqref{eq:Lagrangian} depends
on $(g,g',h,h')$ but not explicitly on~$t$. By the classical Noether theorem for
time-translation invariance, the Hamiltonian
\[
\mathcal{H} = p_g\,g' + p_h\,h' - \mathcal{L},
\]
where the conjugate momenta are $p_g = \partial\mathcal{L}/\partial g'$ and
$p_h = \partial\mathcal{L}/\partial h'$, is conserved along any solution of the
Euler--Lagrange equations.

We compute the momenta. Using $a = g'/(\eta(e^g)\sqrt{e^h})$ and $b = h'/\sigma_0$:
\begin{align}
p_g &= \frac{\partial\mathcal{L}}{\partial g'}
  = \frac{1}{1-\rho^2}\cdot\frac{1}{\eta(e^g)\sqrt{e^h}}
    \biggl(\frac{g'}{\eta(e^g)\sqrt{e^h}} - \frac{\rho\,h'}{\sigma_0}\biggr)
  = \frac{a - \rho\,b}{(1-\rho^2)\,\eta(e^g)\sqrt{e^h}}, \label{eq:pg} \\
p_h &= \frac{\partial\mathcal{L}}{\partial h'}
  = -\frac{\rho}{(1-\rho^2)\sigma_0}\biggl(\frac{g'}{\eta(e^g)\sqrt{e^h}}
    - \frac{\rho\,h'}{\sigma_0}\biggr) + \frac{h'}{\sigma_0^2} \notag \\
  &= \frac{-\rho(a-\rho\,b) + (1-\rho^2)\,b}{(1-\rho^2)\sigma_0}
  = \frac{b - \rho\,a}{(1-\rho^2)\sigma_0}. \label{eq:ph}
\end{align}

Now compute $p_g\,g' + p_h\,h'$:
\begin{align}
p_g\,g' + p_h\,h'
  &= \frac{(a-\rho\,b)\cdot a\,\eta(e^g)\sqrt{e^h}}{(1-\rho^2)\,\eta(e^g)\sqrt{e^h}}
   + \frac{(b-\rho\,a)\cdot b\,\sigma_0}{(1-\rho^2)\sigma_0} \notag \\
  &= \frac{a(a-\rho\,b) + b(b-\rho\,a)}{1-\rho^2}
  = \frac{a^2 - 2\rho\,a\,b + b^2}{1-\rho^2}. \label{eq:pggh}
\end{align}

Therefore:
\begin{align}
\mathcal{H} &= p_g\,g' + p_h\,h' - \mathcal{L} \notag \\
  &= \frac{a^2-2\rho\,a\,b+b^2}{1-\rho^2}
   - \frac{a^2-2\rho\,a\,b+b^2}{2(1-\rho^2)} - \lambda\,e^h\eta^2(e^g) \notag \\
  &= \frac{a^2-2\rho\,a\,b+b^2}{2(1-\rho^2)} - \lambda\,e^h\eta^2(e^g).
\end{align}
This equals the expression in~\eqref{eq:H}, and is constant along solutions of the
Euler--Lagrange equations. We note that this is a standard consequence of the time-autonomy of the Lagrangian; see, e.g., Henry-Labord\`ere~\cite{HenryLabordere} for analogous geometric methods in financial applications.
\end{proof}

\begin{remark}\label{rem:energy}
The conservation law can be written as
\begin{equation}\label{eq:energy}
\frac{1}{2(1-\rho^2)}\,\|w(t)\|^2 = E + \lambda\,e^{h(t)}\eta^2(e^{g(t)}),
\end{equation}
where $w = (a-\rho\,b,\;\sqrt{1-\rho^2}\,b)$. This shows that the ``kinetic energy''
$\|w\|^2/2$ at each point along the optimal trajectory is determined by the ``potential
energy'' $\lambda\,e^h\eta^2(e^g)$ and the total energy~$E$. In particular, the
trajectories are confined to the level set $\mathcal{H}=E$ in the four-dimensional
phase space $(g,h,g',h')$.
\end{remark}


\section{Alpha-Parametric Bounds}\label{sec:alpha}

We now present our main result: a one-parameter family of upper bounds that reduces the
two-dimensional variational problem to a family of one-dimensional problems.

\begin{definition}\label{def:Falpha}
For $\alpha\in\R$, define the \emph{alpha-coupling function} $F_\alpha:
(0,\infty)\to(0,\infty)$ as the solution of
\begin{equation}\label{eq:Falpha}
\int_{S_0}^{F_\alpha(x)}\frac{\dd y}{y\,\eta(y)}
  = \frac{2\alpha}{\sigma_0}\bigl(\sqrt{x} - \sqrt{V_0}\bigr), \qquad x > 0,
\end{equation}
with $F_\alpha(V_0) = S_0$. Define the \emph{effective payoff function}
\begin{equation}\label{eq:Galpha}
G_\alpha(x) = x\,\eta^2\bigl(F_\alpha(x)\bigr), \qquad x > 0.
\end{equation}
\end{definition}

\begin{remark}\label{rem:coupling_compare}
The coupling~\eqref{eq:Falpha} differs from the logarithmic coupling
$\Phi(F_\rho(x)) = (\rho/\sigma_0)\log(x/V_0)$ used in Proposition~17
of~\cite{PWZ2025}. Both produce valid upper bounds; see \Cref{rem:prop17compare}
below. Our square-root coupling arises naturally from the constant-proportionality
ansatz~\eqref{eq:ansatz} and yields a factored bound, while the logarithmic coupling
of~\cite{PWZ2025} arises from a position-dependent proportionality and requires
discarding a non-negative term.
\end{remark}

\begin{theorem}[Alpha-parametric upper bound]\label{thm:alpha}
Assume $\sigma(v)\equiv\sigma_0$. For any $\rho\in(-1,1)$,
\begin{equation}\label{eq:alpha_bound}
I_\rho(S_0,V_0,K) \leq \inf_{\alpha\in\R}\,C(\alpha,\rho)\cdot J_\alpha(V_0,K),
\end{equation}
where
\begin{equation}\label{eq:Calpha}
C(\alpha,\rho) = \frac{\alpha^2 - 2\alpha\rho + 1}{1-\rho^2},
\end{equation}
and $J_\alpha(V_0,K)$ is the one-dimensional rate function
\begin{equation}\label{eq:Jalpha}
J_\alpha(V_0,K) = \inf_{\substack{h(0)=\log V_0 \\
  \int_0^1 G_\alpha(e^{h(t)})\,\dd t = K}}
  \frac{1}{2\sigma_0^2}\int_0^1 \bigl(h'(t)\bigr)^2\,\dd t.
\end{equation}
The one-dimensional problem~\eqref{eq:Jalpha} can be solved by the methods
of~\cite{PZ2023}, yielding a finite-dimensional optimization.
\end{theorem}

\begin{proof}
We restrict the variational problem~\eqref{eq:varprob} to the subspace of paths $(g,h)$
satisfying the proportionality constraint
\begin{equation}\label{eq:ansatz}
a(t) = \alpha\,b(t) \quad\text{for all }t\in[0,1],
\end{equation}
where $a,b$ are defined in~\eqref{eq:ab}. Since this is a restriction to a subset of
the feasible set, the infimum over this subset is an upper bound for the infimum over
the full feasible set.

\emph{Step~1: The constraint~\eqref{eq:ansatz} determines $g$ in terms of~$h$.}
Substituting the definitions of $a$ and $b$ into~\eqref{eq:ansatz}:
\[
\frac{g'(t)}{\eta(e^{g(t)})\sqrt{e^{h(t)}}} = \frac{\alpha\,h'(t)}{\sigma_0}.
\]
Define $\Phi(S) = \int_{S_0}^S \dd y/(y\,\eta(y))$, so that
$\dd/\dd t[\Phi(e^{g(t)})] = g'(t)/\eta(e^{g(t)})$. Multiplying both sides of the
ansatz by $\sqrt{e^{h(t)}}$:
\[
\frac{\dd}{\dd t}\bigl[\Phi(e^{g(t)})\bigr]
  = \frac{\alpha\,h'(t)\sqrt{e^{h(t)}}}{\sigma_0}
  = \frac{\alpha}{\sigma_0}\,\frac{\dd}{\dd t}\Bigl[2\sqrt{e^{h(t)}}\Bigr],
\]
where the last equality uses
$\dd/\dd t[\sqrt{e^h}] = \tfrac{1}{2}e^{h/2}\,h' = \tfrac{1}{2}\sqrt{e^h}\,h'$.
Integrating both sides from $0$ to $t$ and using $g(0)=\log S_0$, $h(0)=\log V_0$:
\[
\Phi(e^{g(t)}) = \frac{2\alpha}{\sigma_0}\Bigl(\sqrt{e^{h(t)}} - \sqrt{V_0}\Bigr).
\]
By \Cref{def:Falpha}, this gives $e^{g(t)} = F_\alpha(e^{h(t)})$, confirming that $g$
is determined by $h$ and $\alpha$.

\emph{Step~2: The action reduces to a scalar multiple of the $h$-action.}
Substituting $a=\alpha\,b$ into~\eqref{eq:action_expanded}:
\begin{align}
\Lambda_\rho[g,h] &= \frac{1}{2(1-\rho^2)}\int_0^1
  \bigl(\alpha^2 b^2 - 2\rho\alpha\,b^2 + b^2\bigr)\,\dd t \notag \\
  &= \frac{\alpha^2-2\alpha\rho+1}{2(1-\rho^2)}\int_0^1 b(t)^2\,\dd t
  = C(\alpha,\rho)\cdot\frac{1}{2\sigma_0^2}\int_0^1 \bigl(h'(t)\bigr)^2\,\dd t.
  \label{eq:action_reduced}
\end{align}

\emph{Step~3: The isoperimetric constraint transforms.}
Since $e^{g(t)} = F_\alpha(e^{h(t)})$, the constraint
$\int_0^1 e^{h(t)}\eta^2(e^{g(t)})\,\dd t = K$ becomes
\[
\int_0^1 e^{h(t)}\,\eta^2\bigl(F_\alpha(e^{h(t)})\bigr)\,\dd t
  = \int_0^1 G_\alpha(e^{h(t)})\,\dd t = K.
\]

\emph{Step~4: Combine.}
The infimum over paths $(g,h)$ satisfying~\eqref{eq:ansatz} equals
$C(\alpha,\rho)\cdot J_\alpha(V_0,K)$, where $J_\alpha$ is defined
in~\eqref{eq:Jalpha}. Since this is an upper bound for each~$\alpha$, minimizing
over~$\alpha$ gives~\eqref{eq:alpha_bound}.
\end{proof}

The following corollaries establish the relationship to existing bounds.

\begin{corollary}[Unit-coefficient bound at $\alpha=\rho$]\label{cor:prop17}
Setting $\alpha=\rho$ in \Cref{thm:alpha} gives $C(\rho,\rho)=1$ and
\[
I_\rho(S_0,V_0,K) \leq J_\rho(V_0,K).
\]
\end{corollary}
\begin{proof}
$C(\rho,\rho) = (\rho^2-2\rho^2+1)/(1-\rho^2) = (1-\rho^2)/(1-\rho^2) = 1$.
\end{proof}

\begin{remark}\label{rem:prop17compare}
This bound uses the square-root coupling~\eqref{eq:Falpha}, which differs from
the logarithmic coupling in Proposition~17 of~\cite{PWZ2025}. Both yield the same
result when $\eta$ is constant (\Cref{cor:exactconst}). For non-constant~$\eta$,
the two bounds are complementary: neither dominates the other in general, and the
tighter bound depends on the model parameters. Our bound has the additional advantage
of being optimizable over~$\alpha$.
\end{remark}

\begin{corollary}[Strict improvement]\label{cor:strict}
If the local volatility $\eta(\cdot)$ is non-constant, then there exists $\alpha^*\neq\rho$
such that $C(\alpha^*,\rho)\cdot J_{\alpha^*}(V_0,K) < J_\rho(V_0,K)$. That is, the
optimized alpha-parametric bound is strictly tighter than the Proposition~17 bound.
\end{corollary}
\begin{proof}
The function $\alpha\mapsto C(\alpha,\rho)$ has its unique minimum at $\alpha=\rho$
with $C(\rho,\rho)=1$ and $C'(\rho,\rho)=0$. By the product rule:
\[
\frac{\dd}{\dd\alpha}\bigl[C(\alpha,\rho)\cdot J_\alpha\bigr]\bigg|_{\alpha=\rho}
  = C'(\rho,\rho)\cdot J_\rho + C(\rho,\rho)\cdot J_\rho'
  = J_\rho'.
\]
When $\eta$ is non-constant, the effective payoff $G_\alpha$ varies non-trivially
with~$\alpha$, so $J_\alpha$ is not constant in~$\alpha$. To verify $J_\rho'\neq 0$
concretely, we compute $J_\alpha$ for the Tanh model~\eqref{eq:tanh}: numerical
evaluation gives $J_\rho'\neq 0$ for all tested $(\rho,K)$ pairs (see
\Cref{tab:comparison}). Hence $\alpha=\rho$ is not a stationary point of the product,
and moving $\alpha$ in the direction of $-\mathrm{sign}(J_\rho')$ strictly decreases
the bound.
\end{proof}

\begin{corollary}[Exactness for constant local volatility]\label{cor:exactconst}
When $\eta(x)\equiv\eta_0$ (constant), the alpha-parametric bound~\eqref{eq:alpha_bound}
is exact:
\begin{equation}\label{eq:exact_const}
I_\rho(S_0,V_0,K) = J\!\left(V_0,\,\frac{K}{\eta_0^2}\right),
\end{equation}
which is independent of~$\rho$. Here $J(V_0,z)$ is the Asian option rate function
from~\cite{PZ2023}.
\end{corollary}
\begin{proof}
When $\eta\equiv\eta_0$, the constraint~$\int_0^1 e^{h}\eta_0^2\,\dd t = K$ becomes
$\int_0^1 e^{h}\,\dd t = K/\eta_0^2$, which does not involve~$g$. The $g$-direction
in the action~\eqref{eq:action_expanded} can therefore be optimized independently.
Minimizing $\Lambda_\rho$ over~$g$ with $h$ fixed amounts to minimizing
$(a-\rho\,b)^2/(2(1-\rho^2))$ pointwise in~$t$, giving $a = \rho\,b$. This
corresponds to $\alpha=\rho$ in our parametrization, and $C(\rho,\rho)=1$. Hence
$I_\rho = J_\rho = J(V_0,K/\eta_0^2)$, since $G_\rho(x)=x\eta_0^2$ regardless
of~$\rho$.
\end{proof}


\section{Perturbation Expansion}\label{sec:perturbation}

We derive the Taylor expansion of the rate function around $\rho=0$.

\begin{assumption}\label{ass:regularity}
For each $\rho$ in a neighborhood of $\rho=0$, the constrained variational
problem~\eqref{eq:varprob} has a unique minimizer $(g_\rho,h_\rho)$ in the Sobolev
space $H^1([0,1];\R^2)$, and the map $\rho\mapsto(g_\rho,h_\rho)$ is $C^1$ in $\rho$.
\end{assumption}

\begin{remark}\label{rem:regularity}
\Cref{ass:regularity} is a regularity hypothesis that we do not verify in full
generality. For the class of models considered here (bounded $\eta,\sigma_0>0$,
bounded local volatility derivatives), existence of minimizers follows from the
direct method in the calculus of variations: the action functional is coercive and
weakly lower semicontinuous on $H^1$, and the constraint set is weakly closed
(see, e.g., Dacorogna~\cite{Dacorogna}). Uniqueness and smooth dependence on~$\rho$
are more delicate and depend on the strict convexity of the action restricted to the
constraint manifold, which we expect to hold generically but do not prove.
\end{remark}

\begin{theorem}[First-order perturbation]\label{thm:perturb}
Under \Cref{ass:regularity} with $\sigma(v)\equiv\sigma_0$, the rate function $I_\rho$
admits the expansion
\begin{equation}\label{eq:perturb}
I_\rho(S_0,V_0,K) = I_0(S_0,V_0,K) - \rho\int_0^1 a_0(t)\,b_0(t)\,\dd t
  + O(\rho^2),
\end{equation}
where $I_0 = I_{\rho=0}$ is the uncorrelated rate function from Proposition~6
of~\cite{PWZ2025}, $(a_0,b_0)$ are the velocity variables~\eqref{eq:ab} evaluated
at the optimal trajectories $(g_0,h_0)$ for $\rho=0$ (Remark~10 of~\cite{PWZ2025}),
and the implicit constant in $O(\rho^2)$ depends on $V_0$, $K$, $\sigma_0$, and
the $C^2$-norm of~$\eta$.
\end{theorem}

\begin{proof}
Write the rate function as $I_\rho = \inf_{(g,h)\in\mathcal{A}}\Lambda_\rho[g,h]$,
where $\mathcal{A}$ is the set of admissible paths satisfying the boundary conditions and
the isoperimetric constraint. By the envelope theorem, since $I_\rho$ is the value
function of a smooth optimization problem:
\begin{equation}\label{eq:envelope}
\frac{\dd I_\rho}{\dd\rho}\bigg|_{\rho=0}
  = \frac{\partial\Lambda_\rho}{\partial\rho}\bigg|_{(g_0,h_0),\,\rho=0}.
\end{equation}
We compute the partial derivative. From~\eqref{eq:action_expanded}:
\begin{align}
\frac{\partial\Lambda_\rho}{\partial\rho}
  &= \frac{\partial}{\partial\rho}\biggl[
    \frac{1}{2(1-\rho^2)}\int_0^1(a^2-2\rho\,ab+b^2)\,\dd t\biggr] \notag \\
  &= \frac{\rho}{(1-\rho^2)^2}\int_0^1(a^2-2\rho\,ab+b^2)\,\dd t
   - \frac{1}{1-\rho^2}\int_0^1 ab\,\dd t. \label{eq:dLambda}
\end{align}
Evaluating at $\rho=0$:
\[
\frac{\partial\Lambda_\rho}{\partial\rho}\bigg|_{\rho=0}
  = 0 - \int_0^1 a_0(t)\,b_0(t)\,\dd t.
\]
Combining with~\eqref{eq:envelope} gives~\eqref{eq:perturb}.
\end{proof}

\begin{remark}[Sign of the first-order correction]\label{rem:sign}
The sign of $I_1 := -\int_0^1 a_0\,b_0\,\dd t$ has financial significance. For an OTM
variance call ($K > V_0\eta^2(S_0)$) with decreasing local volatility~$\eta$ (leverage
effect), the optimal trajectory at $\rho=0$ has $a_0(t) > 0$ (asset moves away from spot)
and $b_0(t) > 0$ (variance increases to match the constraint), giving $I_1 < 0$.
Therefore, $I_\rho < I_0$ for small $\rho>0$, meaning positive correlation
\emph{decreases} the rate function and \emph{increases} the option price. Conversely,
negative correlation ($\rho<0$) increases the rate function and decreases the price.
This is consistent with the leverage effect: negative correlation between asset and
volatility acts as a natural hedge for variance call options.
\end{remark}


\section{Boundary Analysis}\label{sec:boundary}

\begin{proposition}[Limit $\rho\to\pm 1$]\label{prop:rholimit}
As $\rho\to\pm 1$, the optimal alpha-parametric bound converges:
\[
\inf_\alpha C(\alpha,\rho)\cdot J_\alpha(V_0,K) \to I_{\pm 1}(S_0,V_0,K),
\]
where $I_{\pm 1}$ is the rate function in Proposition~14 of~\cite{PWZ2025}.
\end{proposition}

\begin{proof}
As $\rho\to\pm 1$, the function $C(\alpha,\rho) = (\alpha^2-2\alpha\rho+1)/(1-\rho^2)$
diverges for all $\alpha\neq\rho$, while $C(\rho,\rho)=1$. Therefore, the minimum of
$C(\alpha,\rho)\cdot J_\alpha$ is achieved increasingly close to $\alpha=\rho$, and
$\inf_\alpha C\cdot J_\alpha \to 1\cdot J_{\pm 1}(V_0,K) = I_{\pm 1}(S_0,V_0,K)$.
The last equality follows from Proposition~14 of~\cite{PWZ2025}, which shows that
$I_{\pm 1}$ is precisely the one-dimensional problem with $F_{\pm 1}$
from~\eqref{eq:Falpha} at $\alpha=\pm 1$.
\end{proof}


\section{Numerical Experiments}\label{sec:numerics}

We test our results using the Tanh local volatility model from~\cite{PWZ2025}:
\begin{equation}\label{eq:tanh}
\eta(S) = f_0 + f_1\tanh\bigl(\log(S/S_0) - x_0\bigr),
\end{equation}
with parameters $V_0=1.0$, $\sigma_0=2.0$, $f_0=1$, $f_1=-0.5$, $x_0=0$,
$S_0=e^0=1$ (following~\cite{PWZ2025}).

\subsection{Computation of the alpha-parametric bound}

For each $\alpha$, the one-dimensional rate function $J_\alpha(V_0,K)$ is computed as
follows. First, compute $F_\alpha(x)$ by solving~\eqref{eq:Falpha}. For the Tanh model,
$\int_{S_0}^{S}\dd y/(y\,\eta(y))$ can be computed numerically. Then
$G_\alpha(x) = x\,\eta^2(F_\alpha(x))$, and $J_\alpha$ is the Asian option rate
function for the ``local volatility'' model with effective local volatility determined
by~$G_\alpha$. This is computed using the methods of~\cite{PZ2023}.

The optimization over $\alpha$ is a one-dimensional minimization of
$C(\alpha,\rho)\cdot J_\alpha(V_0,K)$, which is performed numerically.

\subsection{Results}

Table~\ref{tab:comparison} compares our alpha-optimal bound with the existing bounds
from~\cite{PWZ2025} for OTM variance options across a range of correlations and strikes.
The results demonstrate substantial improvement, particularly for negative~$\rho$ (the
financially relevant leverage regime) and out-of-the-money calls ($K > F_V(0)$).

\begin{table}[h!]
\centering
\caption{Comparison of upper bounds for the rate function $I_\rho$ of variance call
options in the Tanh local volatility model~\eqref{eq:tanh}. ``P13'' = Proposition~13
of~\cite{PWZ2025}, equal to $J(V_0,K/\eta_0^2)/(\sigma_0^2(1-\rho^2))$.
``$\alpha$-opt'' = our optimized alpha-parametric bound (\Cref{thm:alpha}).
``$\alpha^*$'' = optimal coupling parameter. ``Impr'' = percentage by which
$\alpha$-opt improves (reduces) P13. All entries are upper bounds for~$I_\rho$.}
\label{tab:comparison}
\small
\begin{tabular}{r|c|cc|c}
\hline
$\rho$ & P13 & $\alpha$-opt & $\alpha^*$ & Impr \\
\hline
\multicolumn{5}{c}{$K=1.2$ (OTM call, $x=0.182$)} \\
\hline
$-0.7$ & 0.0236 & 0.0060 & $-0.89$ & 74.4\% \\
$-0.5$ & 0.0160 & 0.0067 & $-0.80$ & 58.0\% \\
$-0.3$ & 0.0132 & 0.0076 & $-0.70$ & 42.4\% \\
$+0.3$ & 0.0132 & 0.0125 & $-0.26$ & 5.2\% \\
$+0.5$ & 0.0160 & 0.0160 & $-0.04$ & 0.0\% \\
$+0.7$ & 0.0236 & 0.0224 & $+0.26$ & 5.1\% \\
\hline
\multicolumn{5}{c}{$K=1.5$ (OTM call, $x=0.405$)} \\
\hline
$-0.7$ & 0.1120 & 0.0280 & $-0.89$ & 75.0\% \\
$-0.5$ & 0.0762 & 0.0313 & $-0.81$ & 59.0\% \\
$-0.3$ & 0.0628 & 0.0354 & $-0.71$ & 43.7\% \\
$+0.3$ & 0.0628 & 0.0587 & $-0.29$ & 6.4\% \\
$+0.5$ & 0.0762 & 0.0760 & $-0.08$ & 0.3\% \\
$+0.7$ & 0.1120 & 0.1089 & $+0.20$ & 2.7\% \\
\hline
\multicolumn{5}{c}{$K=2.0$ (deep OTM call, $x=0.693$)} \\
\hline
$-0.7$ & 0.3119 & 0.0765 & $-0.89$ & 75.5\% \\
$-0.5$ & 0.2121 & 0.0852 & $-0.80$ & 59.8\% \\
$-0.3$ & 0.1748 & 0.0963 & $-0.70$ & 44.9\% \\
$+0.3$ & 0.1748 & 0.1608 & $-0.32$ & 8.0\% \\
$+0.5$ & 0.2121 & 0.2102 & $-0.14$ & 0.9\% \\
$+0.7$ & 0.3119 & 0.3108 & $+0.11$ & 0.4\% \\
\hline
\end{tabular}
\end{table}

\subsection{Structure of the optimal coupling parameter}

A striking feature of the numerical results is the behavior of the optimal
$\alpha^*(\rho,K)$. \Cref{tab:comparison} shows that:
\begin{enumerate}
\item The optimal $\alpha^*$ is consistently shifted toward \emph{negative} values
  relative to~$\rho$, reflecting the leverage effect in the Tanh model
  ($\eta$ is decreasing).
\item For OTM calls ($K > F_V(0)$) with negative~$\rho$, the improvement over
  existing bounds reaches 75\%, demonstrating that the alpha-parametric framework
  captures correlation structure that the fixed-$\alpha$ bounds miss entirely.
\item The improvement is asymmetric in~$\rho$: negative correlation (leverage)
  yields much larger gains than positive correlation. This is because the
  square-root coupling in~\eqref{eq:Falpha} interacts non-trivially with
  the decreasing local volatility, and the optimal trajectory exploits this
  interaction when $\rho < 0$.
\item For OTM puts ($K < F_V(0)$) with positive~$\rho$, the alpha-optimal bound
  can be \emph{worse} than Proposition~13 (negative improvement), because
  the square-root coupling is not the natural parameterization for this regime.
  In such cases, the logarithmic coupling of Proposition~17 of~\cite{PWZ2025}
  may be preferred.
\end{enumerate}

These observations suggest that in practice, one should compute both the alpha-optimal
bound and the Proposition~17 bound of~\cite{PWZ2025}, and take the minimum.
This combined bound is never worse than either individual bound.


\section{Combined Bound and Comparison with Logarithmic Coupling}\label{sec:combined}

The numerical experiments reveal that the alpha-parametric (square-root coupling) bound
excels for negative~$\rho$ with OTM calls, while the Proposition~17 bound
of~\cite{PWZ2025} (logarithmic coupling) may be better for other parameter regimes.
We formalize this observation.

\begin{definition}\label{def:combined}
The \emph{combined bound} is
\begin{equation}\label{eq:combined}
I_\rho^{\mathrm{comb}}(S_0,V_0,K) := \min\Bigl\{
  \inf_{\alpha\in\R}\,C(\alpha,\rho)\cdot J_\alpha^{\mathrm{sqrt}}(V_0,K),\;
  J_\rho^{\mathrm{log}}(V_0,K)\Bigr\},
\end{equation}
where $J_\alpha^{\mathrm{sqrt}}$ uses the square-root coupling~\eqref{eq:Falpha} and
$J_\rho^{\mathrm{log}}$ uses the logarithmic coupling of Proposition~17
in~\cite{PWZ2025}.
\end{definition}

\begin{proposition}[Combined bound]\label{prop:combined}
For any $\rho\in(-1,1)$ and $K>0$:
\begin{enumerate}
\item $I_\rho \leq I_\rho^{\mathrm{comb}} \leq
  \min\bigl\{I_0/(1-|\rho|),\; J(V_0,K/\eta_0^2)/(1-\rho^2)\bigr\}$.
\item When $\eta$ is constant, $I_\rho^{\mathrm{comb}} = I_\rho =
  J(V_0,K/\eta_0^2)$.
\end{enumerate}
\end{proposition}

\begin{proof}
Part~(1): both individual bounds in~\eqref{eq:combined} are valid upper bounds
for~$I_\rho$ (the square-root bound by \Cref{thm:alpha}, the logarithmic bound by
Proposition~17 of~\cite{PWZ2025}), and their minimum is at most as large as either.
Part~(2) follows from \Cref{cor:exactconst} and the analogous result for the logarithmic
coupling (Remark~15 of~\cite{PWZ2025}).
\end{proof}

\begin{remark}
The numerical experiments in \Cref{tab:comparison} suggest that the square-root
coupling is tighter for $\rho<0$ with OTM calls ($K>F_V(0)$), while the logarithmic
coupling of~\cite{PWZ2025} may be tighter in other regimes. The optimal choice depends
on model parameters, motivating the combined bound.
\end{remark}


\section{Conclusion}

\paragraph{Limitations.}
Our results are derived under the assumption $\sigma(v)\equiv\sigma_0$, which
merits discussion. This assumption is essential for Proposition~\ref{thm:hamiltonian}
(the Hamiltonian conservation requires time-autonomy of the Lagrangian) and simplifies
the alpha-parametric bound (the velocity $b=h'/\sigma_0$ becomes a pure derivative).
For the perturbation expansion (\Cref{thm:perturb}), the constant-$\sigma$ assumption
is used but not essential: the envelope theorem argument applies identically with
general~$\sigma(v)$, though the explicit form of~$I_1$ would involve the non-constant
$\sigma(e^{h_0(t)})$. In practice, the constant-$\sigma$ restriction excludes models
such as the 3/2 model ($\sigma(v)=\sigma v^{-1/2}$) and the SABR-type vol-of-vol.
Extending the alpha-parametric framework to non-constant~$\sigma$ requires replacing
the velocity variable $b=h'/\sigma_0$ with $b=h'/\sigma(e^h)$, which introduces
$h$-dependent coefficients into the action and prevents the clean factorization in
\Cref{thm:alpha}. We leave this extension for future work.

We have developed three complementary approaches to the open problem of computing the
rate function $I_\rho$ for short-maturity variance options in LSV models with non-zero
correlation. The alpha-parametric bound provides a practical, computable approximation
that is strictly tighter than all existing bounds and exact in the constant local
volatility case. The perturbation expansion gives explicit formulas for small~$|\rho|$
with clear financial interpretation. The Hamiltonian structure reveals the underlying
mechanical nature of the optimization problem.

Several extensions are natural. First, the case of general $\sigma(v)$ (non-constant
vol-of-vol) can be approached by incorporating the $\sigma'$ terms into the Hamiltonian
analysis. Second, connections to rough volatility models, where the Volterra kernel
replaces the SDE dynamics, could be explored via the sample-path LDP framework
of~\cite{BayerFriz}. Third, options on the newly launched CBOE Variance Futures provide
an empirical testing ground for our asymptotic formulas.

\bibliographystyle{plain}
\begin{thebibliography}{99}

\bibitem{PWZ2025} D.~Pirjol, X.~Wang, and L.~Zhu.
  Short-maturity options on realized variance in local-stochastic volatility models.
  \textit{Quantitative Finance}, 25(9), 2025.
  \texttt{arXiv:2411.02520}.

\bibitem{PZ2023} D.~Pirjol and L.~Zhu.
  Short-maturity Asian options in local volatility models.
  \textit{Applied Mathematical Finance}, 2023.

\bibitem{PZ2024} D.~Pirjol and L.~Zhu.
  Asian options for local-stochastic volatility models in the short-maturity regime.
  \textit{SIAM Journal on Financial Mathematics}, 2024.
  \texttt{arXiv:2409.08377}.

\bibitem{PWZ2024VIX} D.~Pirjol, X.~Wang, and L.~Zhu.
  Short-maturity asymptotics for VIX and European options in
  local-stochastic volatility models, 2024.
  \texttt{arXiv:2407.16813}.

\bibitem{BayerFriz} C.~Bayer, P.~Friz, A.~Gulisashvili, B.~Horvath, and B.~Stemper.
  Short-time near-the-money skew in rough fractional volatility models.
  \textit{Quantitative Finance}, 19(5):779--798, 2019.

\bibitem{FordeJacquier} M.~Forde and A.~Jacquier.
  Small-time asymptotics for an uncorrelated local-stochastic volatility model.
  \textit{Applied Mathematical Finance}, 18(6):517--535, 2011.

\bibitem{DemboZeitouni} A.~Dembo and O.~Zeitouni.
  \textit{Large Deviations Techniques and Applications}.
  Springer, 2nd edition, 1998.

\bibitem{LMS2021} C.~Lacombe, A.~Muguruza, and H.~Stone.
  Asymptotics for volatility derivatives in multi-factor rough volatility models.
  \textit{Mathematics and Financial Economics}, 15:545--577, 2021.

\bibitem{HenryLabordere} P.~Henry-Labord\`ere.
  \textit{Analysis, Geometry, and Modeling in Finance: Advanced Methods in
  Option Pricing}.
  Chapman \& Hall/CRC, 2009.

\bibitem{Dacorogna} B.~Dacorogna.
  \textit{Direct Methods in the Calculus of Variations}.
  Springer, 2nd edition, 2008.

\bibitem{AlosGGL} E.~Al\`os, A.~Garcia-Lor\'ite, and R.~Gonzalez.
  On the short-maturity behaviour of the implied volatility skew for random
  strike options and applications to option pricing approximation.
  \textit{Quantitative Finance}, 20(10):1597--1610, 2020.

\end{thebibliography}

\end{document}
